\documentclass[11pt,letterpaper]{article}
\def\SuiteDocTitle{Agile Requirements Traceability and Governance}
\def\SuiteDocCode{TR-02}
% Shared professional style: pdfLaTeX compatible, no external font files.
\usepackage[T1]{fontenc}
\usepackage[utf8]{inputenc}
\usepackage{lmodern}
\usepackage{microtype}
\usepackage[margin=0.83in,headheight=16pt,footskip=28pt]{geometry}
\usepackage[table]{xcolor}
\definecolor{navy}{HTML}{12304A}
\definecolor{teal}{HTML}{087F83}
\definecolor{gold}{HTML}{B78633}
\definecolor{blueback}{HTML}{F0F5F8}
\definecolor{tealback}{HTML}{EEF8F6}
\definecolor{amberback}{HTML}{FFF7E8}
\definecolor{grayback}{HTML}{F6F7F8}
\usepackage{array,tabularx,longtable,booktabs}
\usepackage{enumitem}
\usepackage{amssymb}
\usepackage{titlesec}
\usepackage{fancyhdr,lastpage}
\usepackage{tcolorbox}
\tcbuselibrary{breakable}
\usepackage{pgffor}
\usepackage{hyperref}
\hypersetup{colorlinks=true,linkcolor=teal,urlcolor=teal,pdfauthor={Engineering Learning Series},pdftitle={\SuiteDocTitle}}
\setlength{\parindent}{0pt}
\setlength{\parskip}{5.5pt}
\setlist{topsep=3pt}
\setcounter{tocdepth}{2}
\titleformat{\section}{\Large\bfseries\color{navy}}{\thesection}{0.7em}{}
\titleformat{\subsection}{\large\bfseries\color{teal}}{\thesubsection}{0.6em}{}
\titleformat{\subsubsection}{\normalsize\bfseries\color{navy}}{\thesubsubsection}{0.6em}{}
\pagestyle{fancy}
\fancyhf{}
\fancyhead[L]{\scriptsize\color{navy}AGILE REQUIREMENTS ENGINEERING}
\fancyhead[R]{\scriptsize\color{teal}\SuiteDocCode}
\fancyfoot[L]{\scriptsize Independent learning and implementation materials}
\fancyfoot[R]{\scriptsize Page \thepage\ of \pageref{LastPage}}
\renewcommand{\headrulewidth}{0.35pt}
\renewcommand{\footrulewidth}{0pt}
\newcommand{\SuiteHeader}[2]{%
{\color{teal}\small\bfseries ENGINEERING DOCUMENTATION SUITE \hfill #2}\par\vspace{8pt}
{\color{navy}\LARGE\bfseries\SuiteDocTitle}\par\vspace{5pt}
{\color{black!68}\normalsize #1}\par\vspace{8pt}
{\color{gold}\rule{\linewidth}{2pt}}\par\vspace{8pt}}
\newcommand{\WritingLines}[1]{\par\foreach \n in {1,...,#1}{\noindent\rule{\linewidth}{0.35pt}\par\vspace{0.28cm}}}
\newcommand{\FreeSpace}[1]{\begin{tcolorbox}[colback=white,colframe=black!25,height=#1,boxrule=0.5pt]\end{tcolorbox}}
\newcommand{\Field}[1]{\noindent\textbf{#1}\quad\rule{0.42\linewidth}{0.35pt}\par\vspace{0.3cm}}
\newcommand{\CheckItem}[1]{\noindent$\square$\quad #1\par\vspace{2pt}}
\newcommand{\SmallNote}[1]{{\footnotesize\color{black!62}#1}\par}
\newcommand{\Guidance}[1]{\begin{tcolorbox}[colback=blueback,colframe=teal,breakable,title=Guidance]#1\end{tcolorbox}}
\sloppy

\begin{document}
\SuiteHeader{Control design | lifecycle evidence and change impact}{Reference}
\tableofcontents\medskip\noindent\rule{\linewidth}{0.4pt}\medskip
\section{Requirements governance intent}
Lightweight agile requirements and durable accountability are compatible. The repository or work tracker acts as an index of evolving intent and supporting evidence. Prefer the smallest useful retained evidence and expand the record when assurance, scale, safety, contractual obligations or team distribution demand it.

\section{Traceability relation}
\begin{longtable}{|p{0.20\linewidth}|p{0.29\linewidth}|p{0.37\linewidth}|}\hline
\textbf{Source} & \textbf{Destination} & \textbf{Trace question} \\ \hline\endfirsthead\hline
\textbf{Source} & \textbf{Destination} & \textbf{Trace question} \\ \hline\endhead
Business outcome & User goal & Which user outcome advances the business purpose? \\ \hline
User goal & Epic / story & Which increments enable the goal? \\ \hline
Story & Acceptance example & How is behavior objectively demonstrated? \\ \hline
Constraint & Test / control & How is the quality obligation verified? \\ \hline
Story & Code / pull request & Which change implements the behavior? \\ \hline
Decision record & Story / architecture & Which tradeoff constrained implementation? \\ \hline
Release & Build, results, approval & What evidence supports delivery? \\ \hline
\end{longtable}
\section{Change-impact protocol}
\begin{enumerate}[leftmargin=*,itemsep=3pt]
\item Log the request and affected user goal; distinguish defect, enhancement and new obligation.
\item Find linked stories, constraints, interfaces, tests and stakeholder decisions.
\item Ask whether the change invalidates prior acceptance or security evidence.
\item Resize and reorder impacted items with the product decision owner.
\item Update acceptance tests and permanent requirement records before release.
\item Retain a dated decision trail and close superseded links.
\end{enumerate}
\section{Definition of Ready versus Definition of Done}
A readiness checklist is an optional team collaboration heuristic, not an excuse to block discovery indefinitely. Done requires an integrated, demonstrable result plus agreed verification and applicable quality controls. Distinguish customer acceptance from deployment approval and from engineering completion.

\section{Evidence retention matrix}
\begin{longtable}{|p{0.20\linewidth}|p{0.29\linewidth}|p{0.37\linewidth}|}\hline
\textbf{Record} & \textbf{Keep?} & \textbf{Rationale} \\ \hline\endfirsthead\hline
\textbf{Record} & \textbf{Keep?} & \textbf{Rationale} \\ \hline\endhead
Accepted story and identifier & Usually & Explains delivered scope \\ \hline
Final acceptance examples and results & Yes & Regression and release evidence \\ \hline
Cross-cutting requirement & Yes & Durable obligation \\ \hline
Conversation notes & Selective & Retain material decisions; discard noise \\ \hline
Work decomposition tasks & Policy-dependent & Usually transient implementation coordination \\ \hline
Risk exception & Yes until closed per policy & Governance and audit \\ \hline
Unvalidated idea & Time-boxed & Avoid misleading stale backlog \\ \hline
\end{longtable}
\section{Review cadence and metrics}
Review new ideas continuously; refine near-term candidates at least once per iteration; inspect acceptance outcomes at iteration review and review cross-cutting constraints before releases. Useful signals include age of blocked items, reopened stories, acceptance defect escapes, story size distribution, throughput variability and missing requirement-test links. Avoid treating point totals as individual performance metrics.

\section{Modern extensions versus the book}
Traceability link graphs, pull-request controls, security gating and repository metadata patterns in this reference are new engineering recommendations. The 2004 source emphasizes useful conversation, planning and confirmation; it does not mandate modern repository workflows or any particular issue-tracker schema.


\end{document}
